\section{Background and Taxonomy}
% Length: ~10%
% Focus: Define the scope and provide the structural map for the whole paper.

% Must contain:
% 1. Definition: A brief explanation of the standard workflow for automated evaluation in RAG.
% 2. Taxonomy Tree Figure: A mandatory diagram showing the three branches of this survey.
% 3. Navigation: A short paragraph referencing the figure and telling the reader how the rest of the paper is organized.

In a standard Retrieval-Augmented Generation pipeline, automated evaluation typically relies on the LLM-as-a-Judge paradigm. The workflow involves three inputs to the evaluating model: a user query, one or more retrieved reference documents, and the target model's generated response. The judge is then asked to assess the response along multiple dimensions. The two most fundamental dimensions are \textit{faithfulness} and \textit{answer relevance}. Faithfulness measures whether the generated answer is grounded in the retrieved documents, while answer relevance determines if it addresses the user's information need \citep{Brehme_2025, gu2025surveyllmasajudge}. A key distinction between standard LLM evaluation and RAG evaluation lies in the evidence sources. RAG evaluators must reconcile two potentially conflicting streams of information: the retrieved context and their own pre-trained parametric knowledge. General-purpose LLM judges already exhibit documented vulnerabilities like verbosity and self-enhancement biases \citep{wang2023fair}. Applying this paradigm to RAG environments introduces further domain-specific vulnerabilities. Existing surveys have yet to categorize these systematically.

Specifically, we identify three interrelated challenges that undermine evaluation reliability in RAG systems. First, long contexts and noisy retrieval outputs push positional bias and distractor sensitivity to the surface. Second, when a judge's internal knowledge contradicts the retrieved evidence, it struggles to decide whether to be truthful or faithful. This ambiguity leads to inconsistent scores. Third, prompt engineering and fine-tuning alone cannot resolve these issues. More capable systems need to query external sources and verify claims actively.

Building on these challenges, we organize the surveyed literature into a three-branch taxonomy, illustrated in Figure~\ref{fig:taxonomy_tree}. This taxonomy maps the complete evaluation lifecycle from the environmental conditions that trigger failures to the underlying mechanisms and the corresponding mitigation strategies.

The remainder of this survey follows this structural map:
% --- Taxonomy Tree Code Start ---
% 如果是双栏论文，请保留 figure* 和 [t]；如果是单栏，请改为 \begin{figure}[htbp]
\begin{figure*}[t]
    \centering
    \resizebox{\textwidth}{!}{%
    \begin{forest}
      for tree={
        grow'=0, draw, thick, rounded corners=3pt, align=center, 
        child anchor=west, parent anchor=east, anchor=west,
        edge path={
          \noexpand\path[\forestoption{edge}, -, thick, draw=gray!80]
            (!u.parent anchor) -- +(4mm,0) |- (.child anchor)\forestoption{edge label};
        },
        % 【修改点】：减小了水平间距和垂直间距，让图片整体变得更紧凑，降低高度
        l sep=5.5mm,  % 原来是 8mm
        s sep=1.5mm,  % 原来是 2.5mm
      },
      leaf style/.style={
        align=left
      }
      [Reliability of\\LLM-as-a-Judge\\in RAG Systems, fill=gray!15, font=\bfseries
        % --- 第一分支：环境诱因 (蓝色系) ---
        [Environmental Triggers\\of Bias, fill=blue!10
          [Context Length \&\\Positional Bias, fill=blue!5
            [Lost in the Middle: \citep{lostinthemiddle2023}, fill=blue!5, leaf style]
            [Distance \& Positional Bias: \citep{distance2024, ragqa2024, posbias2025}, fill=blue!5, leaf style]
            [Attention Calibration: \citep{foundinmiddle2024}, fill=blue!5, leaf style]
          ]
          [Distractors{,} Noise\\\& Cognitive Overload, fill=blue!5
            [Lost in the Noise: \citep{lostinnoise2026}, fill=blue!5, leaf style]
            [Internal Representation Distortion: \citep{internalrep2026}, fill=blue!5, leaf style]
            [Dynamic Context Selection: \citep{dynamiccontext2025}, fill=blue!5, leaf style]
          ]
        ]
        % --- 第二分支：认知冲突 (红色系) ---
        [Epistemological Conflict\\(Truth vs. Faith), fill=red!10
          [Defining the Paradox, fill=red!5
            [Uncertainty Quantification: \citep{fadeeva2026faithfulnessawareuncertaintyquantificationfactchecking}, fill=red!5, leaf style]
            [Conflict Taxonomy: \citep{cattan2025draggedconflictsdetectingaddressing}, fill=red!5, leaf style]
            [TruthfulRAG (Knowledge Graphs): \citep{liu2025truthfulragresolvingfactuallevelconflicts}, fill=red!5, leaf style]
          ]
          [Probing the Black Box, fill=red!5
            [Middle Layers Probing: \citep{gao2025probinglatentknowledgeconflict}, fill=red!5, leaf style]
            [Attention vs. MLPs: \citep{zhao2025understanding}, fill=red!5, leaf style]
            [Transparent Activation Flow: \citep{ye2026seeingconflicttransparentknowledge}, fill=red!5, leaf style]
          ]
          [Dynamic Reliance \& Control, fill=red!5
            [Layer Activation Control: \citep{bi2025parametersvscontextfinegrained}, fill=red!5, leaf style]
            [Parametric Reinforcement: \citep{lin2026resistingcontextualinterferencerag, wang2025accommodateknowledgeconflictsretrievalaugmented}, fill=red!5, leaf style]
            [Micro-Act (Self-Reasoning): \citep{huo2025microactmitigatingknowledgeconflict}, fill=red!5, leaf style]
          ]
        ]
        % --- 第三分支：缓解策略 (绿色系) ---
        [Mitigation to\\Agentic Evaluation, fill=green!10
          [Rubric \& Criteria Calibration, fill=green!5
            [AutoCalibrate: \citep{autocalibrate}, fill=green!5, leaf style]
            [Prometheus 2: \citep{prometheus2}, fill=green!5, leaf style]
          ]
          [Structured Prompting, fill=green!5
            [Multi-dimensional \& Step-by-step: \citep{comm_systems}, fill=green!5, leaf style]
          ]
          [Multi-Agent \& Peer-based, fill=green!5
            [PRD (Peer Rank \& Discussion): \citep{prd}, fill=green!5, leaf style]
            [MACD (Cultural Bias Mitigation): \citep{macd}, fill=green!5, leaf style]
          ]
          [Agentic Verification, fill=green!5
            [Agent-as-a-Judge: \citep{agent_judge}, fill=green!5, leaf style]
            [Auto-Eval Judge (Pipeline): \citep{auto_eval_judge}, fill=green!5, leaf style]
          ]
          [Human-in-the-loop, fill=green!5
            [Ensemble \& Human Oversight: \citep{li2025generationjudgmentopportunitieschallenges}, fill=green!5, leaf style]
          ]
        ]
      ]
    \end{forest}
    }
    \caption{Taxonomy of environmental triggers, internal mechanics, and mitigation strategies for LLM-as-a-Judge in RAG systems.}
    \label{fig:taxonomy_tree}
\end{figure*}
% --- Taxonomy Tree Code End ---
\begin{itemize}    
\item \textbf{Section~\ref{sec:triggers}} examines how extreme context lengths and retrieved noise act as environmental triggers for evaluator bias. We analyze how irrelevant distractors and positional effects systematically degrade the judge's scoring accuracy, and we survey structural interventions that provide only partial remedies.    
\item \textbf{Section~\ref{sec:conflict}} provides a mechanistic analysis of the parametric-contextual knowledge conflict within LLM judges. We synthesize findings from representational probing and interpretability research to explain the underlying reasons for evaluator failure rather than merely documenting the outcomes. This section also reviews dynamic control strategies designed to regulate this tension.    
\item \textbf{Section~\ref{sec:mitigation}} surveys the progression of mitigation strategies from rubric calibration and structured prompting to multi-agent debate and agentic verification pipelines. We analyze the limitations of each approach and argue for the necessity of tool-augmented, stepwise evaluation to address current reliability bottlenecks.
\end{itemize}